\section*{Bab \ref{ch:b2:wittig} --- Reaksi Wittig dan Reaksi Pembentuk C=C Lainnya}

\begin{solution}{exo:b2:wittig:1}
\ce{Ph3P + CH3CH2Br -> [Ph3PCH2CH3]+ + Br-} (substitusi bimolekular,
pemanasan dalam pelarut), lalu butillitium dalam eter kering atau THF di
bawah gas inert: \ce{[Ph3PCH2CH3]+ + BuLi -> Ph3P=CHCH3 + BuH + Li+}.
\end{solution}

\begin{solution}{exo:b2:wittig:2}
Ilida \ce{Ph3P=CH2} dan benzaldehida menghasilkan feniletena (stirena),
\ce{C6H5CH=CH2}, dan trifenilfosfina oksida.
\end{solution}

\begin{solution}{exo:b2:wittig:3}
\ce{Ph3P=CHCOOEt} dan \ce{Ph3P=CHCOCH3}: terstabilkan. \ce{Ph3P=CHCH3}: tak
terstabilkan, yang memerlukan butillitium (atau natrium hidrida, natrium
amida). \ce{Ph3P=CHC6H5}: terstabilkan aril, dalam praktik dibuat dengan
alkoksida atau hidroksida, dan menghasilkan campuran geometri.
\end{solution}

\begin{solution}{exo:b2:wittig:4}
Etil (\textit{E})-pent-2-enoat, \ce{CH3CH2CH=CHCOOEt}, bersama natrium dietil fosfat.
\end{solution}

\begin{solution}{exo:b2:wittig:5}
Wittig: ilida \ce{Ph3P=CH2}, yang dibuat dari metiltrifenilfosfonium iodida
dan butillitium, bersama sikloheksanon hanya menghasilkan
metilenasikloheksana, dengan ikatan rangkap di luar cincin. Dehidrasi
1-metilsikloheksan-1-ol melalui karbokation tersier dan melepaskan hidrogen
yang menghasilkan alkena yang lebih tersubstitusi (Zaitsev): terutama
1-metilsikloheksena, isomer yang salah.
\end{solution}

\begin{solution}{exo:b2:wittig:6}
Kedua potongan itu sama: benzaldehida dan ilida benzilidena
\ce{Ph3P=CHC6H5}. Ilida itu hanya terstabilkan aril dan menghasilkan campuran
(\textit{E})/(\textit{Z}). Untuk (\textit{E}) murni, pakailah reaksi HWE
dietil benzilfosfonat dengan benzaldehida.
\end{solution}

\begin{solution}{exo:b2:wittig:7}
Dengan \ce{Ph3P=CHCH3} terbentuk (\textit{Z})-pent-2-ena,
\ce{CH3CH2CH=CHCH3}, terutama isomer (\textit{Z}). Dengan
\ce{Ph3P=CHCOOEt} terbentuk etil (\textit{E})-pent-2-enoat, terutama
isomer (\textit{E}).
\end{solution}

\begin{solution}{exo:b2:wittig:8}
\ce{Ph3P=CH2 + C6H10O -> C7H12 + Ph3PO}. Massa: metilenasikloheksana
\qty{96.173}{g/mol}, trifenilfosfina oksida \qty{278.291}{g/mol}; ekonomi
atom $96.173/(96.173 + 278.291) =
96.173/374.464 \approx 25.7~\%$. Tiga perempat massa pereaksi keluar sebagai fosfina oksida.
% ledger: aw:C, aw:H, aw:O, aw:P
\end{solution}

\begin{solution}{exo:b2:wittig:9}
Fosfor, yang nukleofilik, menggeser bromida:
\begin{equation*}
\ce{P(OEt)3 + BrCH2COOEt -> [(EtO)3PCH2COOEt]+ + Br-}.
\end{equation*}
Bromida
lalu menyerang sebuah karbon gugus etil pada ion fosfonium itu (substitusi
bimolekular), sambil melepaskan bromoetana dan membentuk ikatan \ce{P=O}:
\begin{equation*}
\ce{[(EtO)3PCH2COOEt]+ + Br- -> (EtO)2P(O)CH2COOEt + EtBr}.
\end{equation*}
\end{solution}

\begin{solution}{exo:b2:wittig:10}
\ce{OHC-CH=CH-CHO + 2Ph3P=CHC6H5 -> C6H5CH=CH-CH=CH-CH=CHC6H5 + 2Ph3PO}:
dua ekuivalen ilida benzilidena (dari benziltrifenilfosfonium klorida dan
sebuah basa), satu untuk setiap gugus aldehida. Ikatan-ikatan rangkap yang
baru terbentuk sebagai campuran geometri; isomer serba-(\textit{E}), yang
paling stabil, diisolasi melalui kristalisasi, atau campurannya
diisomerisasi.
\end{solution}

\begin{solution}{exo:b2:wittig:11}
\omterm{def:b2:enolates-aldol:aldol}{Aldol}: benzaldehida dengan propanon berlebih dan natrium hidroksida
(\cref{ch:b2:enolates-aldol}); pereaksinya murah, produk sampingnya air,
produknya (\textit{E}), tetapi kondensasi dapat terjadi dua kali
(dibenzalaseton) jika propanon tidak berlebih. Wittig: benzaldehida dengan
\omterm{def:b2:wittig:ylide}{ilida terstabilkan}
\ce{Ph3P=CHCOCH3}; (\textit{E}), satu produk, tanpa kondensasi ganda, tetapi
ilidanya mahal dan trifenilfosfina oksida harus disingkirkan. Di sini \omterm{def:b2:enolates-aldol:aldol}{aldol}
adalah jalur yang lebih baik.
\end{solution}

\begin{solution}{exo:b2:wittig:12}
Wittig: butanal dengan \ce{Ph3P=CHCH2CH2CH3} (dari 1-bromobutana, lalu
butillitium), bebas garam: terutama (\textit{Z}); produk sampingnya
trifenilfosfina oksida, dan sebagian isomer (\textit{E}). Alkuna: okt-4-una
(dari etuna melalui dua alkilasi dengan 1-bromopropana dan natrium amida)
yang dihidrogenasi dengan katalis paladium yang diracuni
(\cref{prop:b2:alkene-redox:syn-hydrogenation}): hanya (\textit{Z}); produk
sampingnya natrium bromida dan amonia dari alkilasi.
\end{solution}

\begin{solution}{pb:b2:wittig:1}
\textbf{1.} \ce{C23H46}, yaitu \ce{C_nH_{2n}} dengan $n = 23$: satu ikatan
rangkap (atau satu cincin), di sini sebuah
\ce{C=C}.
\textbf{2.} Ikatan rangkap \ce{C9=C10}.
\textbf{3.} Nonanal \ce{CH3(CH2)7CHO} dengan ilida dari 1-halotetradekana;
atau tetradekanal \ce{CH3(CH2)12CHO} dengan ilida dari 1-halononana.
\textbf{4.} 1-Bromotetradekana, \ce{CH3(CH2)13Br}.
\textbf{5.} Substitusi bimolekular oleh trifenilfosfina yang meruah
berlangsung cepat pada karbon primer yang tidak terhalang dan tidak
tersaingi oleh eliminasi.
\textbf{6.} \ce{Ph3P + C14H29Br -> [Ph3PC14H29]+ + Br-}, tetradesiltrifenilfosfonium bromida.
\textbf{7.} Butillitium (atau natrium hidrida, natrium amida); hidroksida
adalah basa yang terlalu lemah untuk ilida tak terstabilkan, dan air yang
dibawanya akan memprotonasi ilida.
\textbf{8.} Tidak, ilida itu hanya membawa gugus alkil: terutama (\textit{Z})
(\cref{prop:b2:wittig:stereo}).
\textbf{9.} Cincin beranggota empat \ce{P-C-C-O}: fosfor terikat pada karbon
yang membawa rantai \ce{C13H27}, karbon itu pada karbon yang membawa rantai
\ce{C8H17}, dan karbon terakhir ini pada oksigen, yang terikat pada fosfor.
\textbf{10.} Butillitium dan ilida bereaksi dengan air dan oksigen.
\textbf{11.} Ikatan rangkap menghubungkan bekas karbon karbonil nonanal (C9
produk) dengan bekas karbon ilida (C10), dan tidak ada ikatan lain yang
bergeser (\cref{prop:b2:wittig:regiospecific}).
\textbf{12.} Karbokation pada C9, yang dapat kehilangan hidrogen dari C8 atau
C10 dan dapat bergeser: campuran trikos-8-ena dan trikos-9-ena, masing-masing
sebagai (\textit{E}) dan (\textit{Z}), terutama (\textit{E}).
\textbf{13.} \ce{Ph3P=CHC13H27 + C8H17CHO -> C8H17CH=CHC13H27 + Ph3PO}.
\textbf{14.} Target \ce{C23H46}: \qty{322.621}{g/mol}; nonanal \ce{C9H18O}: \qty{142.242}{g/mol}; garam
\ce{C32H44BrP}: \qty{539.582}{g/mol}; \ce{Ph3PO}, \ce{C18H15OP}: \qty{278.291}{g/mol}.
\textbf{15.} Ilida \ce{C32H43P} bermassa $539.582 - 80.912 = \qty{458.670}{g/mol}$; bersama nonanal,
\qty{600.912}{g/mol} pereaksi untuk \qty{322.621}{g/mol} produk: $322.621/600.912 \approx
53.7~\%$.
\textbf{16.} Berdasarkan kepolarannya: hidrokarbon nonpolar lolos melalui
kolom silika pendek dengan eluen nonpolar sementara oksidanya tertahan;
atau oksidanya mengkristal dari pelarut nonpolar yang dingin.
\textbf{17.} Ikatan \ce{P=O} yang sangat kuat yang terbentuk melebihi ikatan
\ce{P-C} yang terputus (\cref{prop:b2:wittig:driving}).
\textbf{18.} Reaksi itu memerlukan \omterm{def:b2:wittig:hwe}{fosfonat} dengan gugus penarik elektron
pada karbon yang bereaksi, yang tidak ada pada rantai alkil biasa, dan
reaksi itu menghasilkan alkena (\textit{E}).
\textbf{19.} Trikos-9-una, \ce{CH3(CH2)7C#C(CH2)12CH3}, dengan hidrogen pada
katalis paladium yang diracuni (paladium pada kalsium karbonat dengan garam
timbal dan kuinolina).
\textbf{20.} Dek-1-una, \ce{CH3(CH2)7C#CH}, dideprotonasi oleh natrium amida,
lalu 1-bromotridekana
\ce{CH3(CH2)12Br}: $8 + 2 + 13 = 23$ karbon, dengan ikatan rangkap tiga di antara C9 dan C10.
\textbf{21.} Hidrogenasi lebih lanjut akan menghasilkan trikosana, tanpa
ikatan rangkap; katalis yang diracuni mengikat alkuna jauh lebih kuat
daripada alkena, yang meninggalkan permukaan sebelum adisi kedua.
\textbf{22.} Keduanya memerlukan dua tahap dari potongan-potongan yang
tersedia. Jalur alkuna menghasilkan isomer (\textit{Z}) dengan bersih dan
hanya menyisakan garam; jalur Wittig menghasilkan sebagian isomer
(\textit{E}) dan setumpuk trifenilfosfina oksida.
\textbf{23.} 100~\%: \ce{C23H44 + H2 -> C23H46}, setiap atom berakhir dalam produk.
\textbf{24.} $278.291/322.621 = \boldsymbol{\approx \qty{0.86}{g}}$
trifenilfosfina oksida per gram feromon.
\end{solution}
